Discrete-choice experiments are widely employed for policy and welfare analysis as they are assumed to reflect people's stable, underlying preferences for specific options' attributes. However, constantly calculating optimal tradeoffs is cognitively demanding and may not persist when repeated choices are made in stable environments. Using eye-tracking and computational modeling in a DCE, we show that extended repetition in stable environments fundamentally changes the cognitive architecture of decision-making, where different behavioral strategies are deployed conditionally on how attentional resources are deployed. While some subjects maintain exploratory attention and behave as utility-maximizers, another group offloads attention from attributes, relying on spatial, habit-like heuristics. Critically, we show that choice inertia---the marker of automaticity---is endogenous to visual gaze entropy. By including changes in environmental stability, attention is captured and habitual subjects instantly break their heuristics. Contrary to widely held assumptions, joint computational modeling demonstrates that preferences are dynamically constructed and affected by the level of attentional load. Our results thus bridge cognitive mechanisms with structural econometrics to resolve a foundational problem in discrete-choice modeling.

